% !TEX TS-program = lualatex
% !TEX encoding = UTF-8 Unicode
\documentclass[ngerman,11pt]{article}
\usepackage{fontspec}
\usepackage[ngerman]{babel}
\usepackage[a4paper, total={16cm, 25cm}]{geometry}
\usepackage{amsmath}
\usepackage{amsfonts}
\usepackage{amssymb}
\usepackage{multicol}
\usepackage{upquote}
\usepackage{tcolorbox}
\usepackage{cancel}
\usepackage{enumitem}
\usepackage{url}
\usepackage[colorlinks=true, linkcolor=blue, urlcolor=blue]{hyperref}
\usepackage{listings}
\lstset{
    language=Python,
    basicstyle=\ttfamily,
    keywordstyle=\color{blue},
    stringstyle=\color{red},
    commentstyle=\color{gray},
    breaklines=true
}
\begin{document}
\hrule
\begin{center}
{\large\bf Python: Grundlegende Datentypen}\\[-0mm]
Studienkolleg der TU Berlin\hfill Till Zoppke\\
Informatik \hfill März 2026\\[2mm]
\end{center}
\hrule
\par\bigskip
\noindent
Die \textbf{Python interaktive Konsole}\footnote{\url{https://docs.python.org/3/tutorial/appendix.html\#tut-interac}} ist ein mächtiges Werkzeug zum schnellen Testen von Code und zum Erlernen der Sprache. Sie können die interaktive Konsole starten, indem Sie in Ihrem Terminal oder Ihrer Kommandozeile \texttt{python} oder \texttt{python3} eingeben. Alternativ können Sie eine web-basierte Installation\footnote{\url{https://trinket.io/console}} nutzen.

\begin{itemize}[itemsep=1pt]
    \item Der Prompt \texttt{>>>} zeigt an, dass Python auf Ihre Eingabe wartet.
    \item Jeder eingegebene Befehl wird sofort ausgeführt.
    \item Für eine spätere Verwendung können Sie Werte in Variablen speichern, z.B. \texttt{n = 3}
    \item Ausdrücke ohne explizite Zuweisung zeigen ihren Wert direkt an.
    \item Der sekundäre Prompt \texttt{...} erscheint bei mehrzeiligen Eingaben.
    \item Verwenden Sie \texttt{help()} für die eingebaute Dokumentation.
    \item Mit \texttt{quit()} oder \texttt{exit()} beenden Sie die Konsole.
\end{itemize}

% Integer
\section{Integer (Ganzzahlen)}

Integer repräsentieren ganze Zahlen, wie wir sie aus der Mathematik kennen. In Python haben Integer keine Größenbeschränkung, sodass auch sehr große Zahlen problemlos verarbeitet werden können. Die Operatoren entsprechen den üblichen arithmetischen Operationen.

\begin{multicols}{2}\raggedcolumns
\noindent\textbf{Arithmetische Operatoren:}
\begin{itemize}[leftmargin=*]
    \item \texttt{+}, \texttt{-} (Vorzeichen)
    \item \texttt{+}, \texttt{-}, \texttt{*} (Grundrechenarten)
    \item \texttt{/} (Division, ergibt Float)
    \item \texttt{//} (Ganzzahlige Division)
    \item \texttt{\%} (Modulo, Rest)
    \item \texttt{**} (Potenzierung)
\end{itemize}
\par\bigskip

\noindent\textbf{Vergleichsoperatoren:}
\begin{itemize}[leftmargin=*]
    \item \texttt{<}, \texttt{>} (kleiner, größer)
    \item \texttt{<=}, \texttt{>=} (kleiner gleich, größer gleich)
    \item \texttt{==}, \texttt{!=} (gleich, ungleich)
\end{itemize}
\par\bigskip

\noindent\textbf{Eingebaute Funktionen:}
\begin{itemize}[leftmargin=*]
    \item \texttt{int()} - Konvertierung zu Integer
    \item \texttt{abs()} - Absolutwert
\end{itemize}
\columnbreak
\begin{lstlisting}
>>> type(42)
<class 'int'>
>>> a = 5
>>> b = -10
>>> a + b, a * b
(-5, -50)
>>> a / 2, a // 2
(2.5, 2)
>>> a % 2
1
>>> a ** 0, a ** 1, a ** 2
(1, 5, 25)
>>> (a + b) * 2
-10
>>> -a
-5
>>> int("42"), abs(b)
(42, 10)
>>> a < b, a == 5
(False, True)
>>> a >= 5 and b <= -10
True
\end{lstlisting}
\end{multicols}

% Boolean
\section{Boolean (Wahrheitswerte)}

Boolean-Werte repräsentieren Wahrheitswerte und können nur zwei Werte annehmen: \texttt{True} oder \texttt{False}. Sie werden für Bedingungen verwendet, beispielsweise bei Verzweigungen oder für Schleifen.

\begin{multicols}{2}\raggedcolumns
\noindent\textbf{Literale:}
\begin{itemize}[leftmargin=*]
    \item \texttt{True, False} - Wahr und Falsch
\end{itemize}
\par\smallskip

\noindent\textbf{Operatoren:}
\begin{itemize}[leftmargin=*]
    \item \texttt{and} - Logisches UND
    \item \texttt{or} - Logisches ODER
    \item \texttt{not} - Logische Negation
\end{itemize}
\par\smallskip

\noindent\textbf{Eingebaute Funktionen:}
\begin{itemize}[leftmargin=*]
    \item \texttt{bool()} - Konvertierung zu Boolean
\end{itemize}
\columnbreak
\begin{lstlisting}
>>> type(True)
<class 'bool'>
>>> x, y = True, False
>>> x and y, x or y
(False, True)
>>> not x, not y
(False, True)
>>> bool(1), bool(0)
(True, False)
>>> bool(""), bool("text")
(False, True)
\end{lstlisting}
\end{multicols}

% String
\section{String (Zeichenketten)}

Strings sind Sequenzen von Zeichen. Jedes einzelne Zeichen wird als Unicode kodiert. Ein String kann beliebig lang sein. Ein Literal kann durch einfache oder doppelte Anführungszeichen begrenzt sein. Mit \texttt{"""} wird ein mehrzeiliger String gekennzeichnet.

\begin{multicols}{2}\raggedcolumns
\noindent\textbf{Literale:}
\begin{itemize}[leftmargin=*]
    \item \texttt{"Bärlin", 'Osolino'}  
    \item \texttt{\char39"\char39, "\char39"} 
\end{itemize}
\par\bigskip

\noindent\textbf{Operatoren:}
\begin{itemize}[leftmargin=*]
    \item \texttt{+} (Verkettung / Konkatenation)
    \item \texttt{*} (Wiederholung)
    \item \texttt{[]} (Indexierung)
    \item \texttt{[:]} (Slicing)
    \item \texttt{in} (Enthaltensein)
\end{itemize}
\par\bigskip

\noindent\textbf{Eingebaute Funktionen:}
\begin{itemize}[leftmargin=*]
    \item \texttt{str()} - Konvertierung zu String
    \item \texttt{len()} - Länge ermitteln
    \item \texttt{ord()}, \texttt{chr()} - Unicode 
\end{itemize}
\columnbreak

\begin{lstlisting}
>>> type("Hallo")
<class 'str'>
>>> s1, s2 = "Hallo", "Welt"
>>> s1 + " " + s2
'Hallo Welt'
>>> 3*'katzenklo'
'katzenklokatzenklokatzenklo'
>>> s1[1], s1[1:4]
('a', 'all')
>>> "a" in s1, "x" in s1
(True, False)
>>> str(42), len(s1)
('42', 5)
>>> ord('A'), chr(65)
(65, 'A')
>>> """
... hier 
... geht es weiter"""
'\nhier \ngeht es weiter'
\end{lstlisting}
\end{multicols}
\pagebreak
% Float
\section{Float (Fließkommazahlen)}

Floats repräsentieren Fließkommazahlen mit begrenzter Genauigkeit. Sie werden für alle Berechnungen verwendet, bei denen Dezimalstellen benötigt werden. Floats in Python entsprechen dem IEEE-754-Standard für Fließkommazahlen\footnote{\url{https://de.wikipedia.org/wiki/IEEE_754}} mit doppelter Genauigkeit. Beim Rechnen kommt es mitunter zu Rundungsfehlern.

\begin{multicols}{2}\raggedcolumns
\noindent\textbf{Operatoren:}
\begin{itemize}[leftmargin=*]
    \item \texttt{+}, \texttt{-}, \texttt{*}, \texttt{/}
    \item \texttt{//} (Division mit Rest)
    \item \texttt{\%} (Modulo)
    \item \texttt{**} (Potenzierung)
    \item Vergleichsoperatoren
\end{itemize}
\par\bigskip
\noindent\textbf{Eingebaute Funktionen:}
\begin{itemize}[leftmargin=*]
    \item \texttt{float()} - Konvertierung zu Float
    \item \texttt{round()} - Runden
    \item \texttt{abs()} - Absolutwert
\end{itemize}
\columnbreak
\begin{lstlisting}
>>> type(0.0)
<class 'float'>
>>> f1, f2 = 3.1415, 2.4
>>> f1 + f2, f1 - f2
(5.5415, 0.7415000000000003)
>>> f1 * f2, f1 / f2
(7.5396, 1.3089583333333334)
>>> f1 // f2
1.0
>>> 7 % 2.4
2.2
>>> 2 ** 0.5
1.4142135623730951
>>> float("-.4"), float('inf')
(-0.4, inf)
>>> round(2 ** 0.5, 3)
1.414
\end{lstlisting}
\end{multicols}

% Tupel
\section{Tuple (Tupel)}

Tupel sind unveränderbare (immutable) Sequenzen, die eine Sammlung von Elementen enthalten. Die Werte schreibt man, durch Kommata getrennt, in runden Klammern. Im Gegensatz zu Listen können Tupel nach ihrer Erstellung nicht mehr verändert werden.

\begin{multicols}{2}\raggedcolumns
\noindent\textbf{Literale:}
\begin{itemize}[leftmargin=*]
    \item \texttt{()} Leeres Tuple
    \item \texttt{(1,)} Tuple mit einem Element
\end{itemize}
\par\bigskip

\noindent\textbf{Operatoren:}
\begin{itemize}[leftmargin=*]
    \item \texttt{+} (Verkettung / Konlatenation)
    \item \texttt{*} (Wiederholung)
    \item \texttt{[]} (Indexierung)
    \item \texttt{[:]} (Slicing)
    \item \texttt{in} (Enthaltensein)
\end{itemize}
\par\bigskip

\noindent\textbf{Eingebaute Methoden:}
\begin{itemize}[leftmargin=*]
    \item \texttt{tuple()} - Konvertierung zu Tuple
    \item \texttt{len()} - Länge ermitteln
\end{itemize}
\columnbreak
\begin{lstlisting}
>>> type((1, 2, 3))
<class 'tuple'>
>>> tup1 = (1, 2, 3)
>>> tup2 = ("a", "b", "c")
>>> tup1 + tup2
(1, 2, 3, 'a', 'b', 'c')
>>> ('x', 'abc', 7) * 2
('x', 'abc', 7, 'x', 'abc', 7)
>>> tup1[0], tup2[0:2]
(1, ('a', 'b'))
>>> 2 in tup1, 5 in tup1
(True, False)
>>> tuple([1, 2, 3]), len(tup1)
((1, 2, 3), 3)
>>> # Unveränderbar:
>>> # tup1[0] = 5  # TypeError
\end{lstlisting}
\end{multicols}

% Listen
\section{List (Listen)}

Listen sind veränderbare (mutable) geordnete Sequenzen. Sie können Elemente unterschiedlicher Datentypen enthalten. Listen sind eine der am häufigsten verwendeten Datenstrukturen in Python und bieten viele Methoden zur Manipulation von Daten.

\begin{multicols}{2}\raggedcolumns
\noindent\textbf{Operatoren:}
\begin{itemize}[leftmargin=*]
    \item \texttt{+} (Verkettung / Konkatenation)
    \item \texttt{*} (Wiederholung)
    \item \texttt{[]} (Indexierung)
    \item \texttt{[:]} (Slicing)
    \item \texttt{in} (Enthaltensein)
\end{itemize}
\par\bigskip

\noindent\textbf{Eingebaute Methoden:}
\begin{itemize}[leftmargin=*]
    \item \texttt{list()} - Konvertierung zu Liste
    \item \texttt{len()} - Länge ermitteln
    \item \texttt{all()} - Prüft ob alle Elemente True sind
    \item \texttt{any()} - Prüft ob mindestens ein Element True ist
    \item \texttt{max()} - Größtes Element
    \item \texttt{min()} - Kleinstes Element
    \item \texttt{sum()} - Summe aller Elemente
\end{itemize}
\columnbreak

\begin{lstlisting}
>>> type([1, 2, 3])
<class 'list'>
>>> a = list(range(1,4))
>>> b = ["a", "b", "c"]
>>> a + b
[1, 2, 3, 'a', 'b', 'c']
>>> a * 2
[1, 2, 3, 1, 2, 3]
>>> a[0], b[0:2]
(1, ['a', 'b'])
>>> 2 in a, 5 in a
(True, False)
>>> list((1, 2, 3)), len(a)
([1, 2, 3], 3)
>>> all([True, True, False])
False
>>> any([False, False, True])
True
>>> max(b), min(b)
(5, 1)
>>> sum(b)
15
>>> bools = [n > 2 for n in b]
>>> bools
[True, True, True, False, False]
\end{lstlisting}
\end{multicols}

\end{document}